% ===========================================================================
% Introduction générale.
% Aucun résultat annoncé ici qui ne soit démontré au chapitre 5.
% ===========================================================================
\chapter*{Introduction générale}
\addcontentsline{toc}{chapter}{Introduction générale}

Le recrutement est devenu une activité à fort volume documentaire. Une offre
d'emploi publiée sur un canal grand public attire couramment plusieurs dizaines
de curriculum vitæ, déposés dans des formats hétérogènes, rédigés selon des
conventions personnelles, et dont la lecture complète mobilise un temps que les
équipes de ressources humaines n'ont pas toujours. Cette réalité crée une
tension bien identifiée : plus le vivier est large, plus la sélection gagne en
qualité potentielle, et moins elle est praticable manuellement.

Les outils existants répondent partiellement à cette tension. Les systèmes de
suivi des candidatures classiques organisent le stockage et le cycle de vie des
dossiers, mais ne lisent pas leur contenu : ils déplacent le problème sans le
traiter. Les filtres par mots-clés introduisent une automatisation, au prix
d'une fragilité connue — un candidat qui écrit « conception d'interfaces » là
où l'annonce demande « UI design » disparaît du résultat, sans que personne ne
s'en aperçoive. Les approches génératives récentes améliorent la souplesse de
lecture, mais produisent des réponses dont rien ne garantit qu'elles reposent
sur le dossier réel du candidat plutôt que sur une plausibilité statistique.

Le problème d'ingénierie se formule alors ainsi : comment traiter
automatiquement un volume de candidatures tout en conservant une décision
explicable, une comparaison reproductible entre candidats, et la possibilité
pour un recruteur de comprendre — donc de contester — chaque résultat produit
par la machine ?

Le projet présenté dans ce rapport traite cette question par la conception et
la réalisation de PIQBIT, une plateforme de recrutement augmentée par
l'apprentissage automatique. PIQBIT ne se limite pas au tri des candidatures :
il couvre la chaîne complète, du dépôt d'un CV jusqu'à l'activation de la fiche
employé, en passant par l'analyse, la notation, le rapprochement sémantique, la
planification d'entretiens, le test technique, la prédiction salariale, la
négociation et la signature électronique du contrat.

L'objectif général de ce projet de fin d'études est de concevoir et de réaliser
une plateforme qui automatise le traitement des candidatures sans sacrifier
l'explicabilité des décisions, et qui accompagne le recruteur sur l'ensemble du
cycle plutôt que sur son seul premier filtre.

Cet objectif se décline en objectifs spécifiques :

\begin{itemize}[leftmargin=*]
  \item bâtir un socle applicatif sécurisé — authentification, second facteur,
        contrôle d'accès par rôle, schéma piloté par migrations ;
  \item mettre en place la gestion des offres, des candidatures et des CV,
        avec validation des dépôts et cycle de vie explicite ;
  \item implémenter l'analyse automatique de CV et une notation
        \emph{explicable}, décomposée en sous-scores rattachés aux exigences
        de l'offre ;
  \item implémenter un rapprochement sémantique vectoriel doté d'un chemin de
        repli fonctionnel ;
  \item couvrir les étapes aval du recrutement : entretiens avec confirmation
        de créneau, test technique, contrats et signature électronique ;
  \item entraîner et intégrer un modèle de prédiction salariale, et en
        justifier la sélection sur des critères mesurés ;
  \item valider la plateforme par des tests automatisés, des portes de qualité
        en intégration continue et une campagne datée.
\end{itemize}

Le rapport s'organise en cinq chapitres. Le \textbf{chapitre~1} pose le
contexte général du recrutement assisté, présente l'organisme d'accueil,
étudie les solutions existantes, formule la problématique et délimite le
périmètre du projet. Le \textbf{chapitre~2} conduit l'ingénierie des
exigences : acteurs, exigences fonctionnelles et non fonctionnelles, cas
d'utilisation et traçabilité. Le \textbf{chapitre~3} expose l'architecture et
la conception : contexte système, vue en composants, modèle de données,
séquences principales, frontières de sécurité et décisions d'architecture. Le
\textbf{chapitre~4} décrit l'implémentation, sous-système par sous-système, en
rapportant les difficultés réellement rencontrées et la manière dont elles ont
été traitées. Le \textbf{chapitre~5} présente la validation, la sécurité et
l'exploitation, avec des résultats datés — y compris ceux qui sont en échec.

La conclusion générale dresse le bilan objectif par objectif, énonce les
limites de la solution livrée et propose des perspectives d'extension.
